% LAYOUT TEST -- not content. Exercises every element of the layout vocabulary so that a change to the preamble
% can be checked at a glance. Remove it from main.tex (and delete it) when the first real chapter exists.
\chapter{Layout test}\label{ch:layout}

This chapter only shows the layout. Running text is set in IBM Plex Sans at 10~pt on a 114~mm block, with
block paragraphs.\sidenote{A numbered sidenote sits in the margin, level with its reference.} Mathematics is
set in a sans-serif face to match:
\begin{equation}
  f'(x) = \frac{f(x+h) - f(x-h)}{2h} - \frac{h^2}{6} f'''(\xi), \qquad \xi \in (x-h, x+h).
  \label{eq:layout-central}
\end{equation}
\mnote{An unnumbered margin note, for a short gloss.}A reference to \cref{eq:layout-central}, to
\cref{sb:layout}, to \cref{fig:layout} and a citation\cite{layout-test} all use the accent colour.

\section{A section}
\subsection{A subsection}
\paragraph{A run-in paragraph heading.} Code is set on a light-gray ground, quoted from a compiled example
(\code{code/quickstart.cpp}) between its snippet markers:
\snippet{quickstart}{quickstart}
A number computed by the library comes from \code{data/values.tex}: the default step gives a relative error of
\valTemplateDefaultError{} here.

\begin{sidebar}[label=sb:layout]{A sidebar on a harder topic}
A sidebar carries a topic that a user does not need, such as why the rounding of $x + h$ matters. The text
around it must read without it. It may hold equations, $\epsmach = 2^{-52}$, and lists:
\begin{itemize}
  \item it is numbered within the chapter and listed in the List of sidebars;
  \item it breaks across pages when it has to.
\end{itemize}
\end{sidebar}

\begin{figure}
\sidecaption{A text-width figure with its caption in the margin: the error of \code{Central<2>} for
$\deriv{}{x}\mathrm{e}^x$ at $x = 1$ against the step, from generated data
(\code{data/template-step-error.csv}).\label{fig:layout}}
\centering
\begin{tikzpicture}
\begin{axis}[eqtufte, width=95mm, height=45mm, xmode=log, ymode=log, xlabel={step $h$}, ylabel={error}]
  \addplot[accent, mark=*] table[col sep=comma, x=h, y=error, comment chars={\#}] {data/template-step-error.csv};
\end{axis}
\end{tikzpicture}
\end{figure}

\begin{margintable}
\caption{A margin table.}
\centering\footnotesize
\begin{tabular}{@{}lr@{}}
\toprule
Rule & $\ampl$\\
\midrule
\code{Central<2>} & 1\\
\code{Central<4>} & 1.5\\
\bottomrule
\end{tabular}
\end{margintable}

\begin{table*}
\caption{A full-width table.}
\centering\small
\begin{tabular}{@{}llll@{}}
\toprule
Method & Nodes & Evaluations & Error estimate\\
\midrule
\code{Central<P>} & $P + N - 1$ & $P + N - 1$ & no\\
\code{Ridders} & $\pm h/2^i$ & 13 to 45 & yes\\
\bottomrule
\end{tabular}
\end{table*}
